\chapter{Reading a Trading Firm's Accounts}\label{ch:in:reading-a-trading-firms-accounts}

For nothing but a download, the UK register gives years of a trading firm's accounts, many of them machine-readable:
turnover, operating profit, the average number of employees, \omterm{def:in:reading-a-trading-firms-accounts:staff}{staff costs}, and for a partnership the profit shared among
its members. In the year to 31 January 2025 one London investment manager, Quadrature Capital Limited, reported \omterm{def:in:reading-a-trading-firms-accounts:staff}{staff
costs} of \pounds556 million for an average of 173 employees: \pounds3.2 million a head. \omterm{def:in:reading-a-trading-firms-accounts:staff}{Staff costs} over average employees is
the most honest pay figure the industry publishes, because the law defines both halves of the fraction. It is also
easy to misread. It is a mean, not what a typical employee earns; it counts employees but not partners; it belongs to
one legal entity, which may not be where the firm's people or profits sit; and one year of a trading firm is not the
next. This chapter shows where the filings are, what they contain, how to read the machine-readable ones with code,
and the pitfalls that turn a correct number into a wrong conclusion. Chapter~12 then puts the whole industry's filed
numbers side by side, and Book~16, chapter~1, is the reference for the economics the lines describe.

\section{Where the filings are: registries, EDGAR, annual reports, prospectuses, divisional and Pillar~3 reports}

\begin{definition}[Company registry, statutory accounts]\label{def:in:reading-a-trading-firms-accounts:registry}
A \emph{company registry}\index{company registry} is the public register with which a country's companies and
partnerships file their constitution, their officers and owners, and their annual accounts. \emph{Statutory
accounts}\index{statutory accounts} are the annual financial statements that the law requires an entity to prepare,
have audited unless it is exempt, and file: a balance sheet, a profit and loss account, notes, and the reports of the
directors (or members) and of the auditor.
\end{definition}

A trading firm is rarely one entity. It is a group: a parent, often in the country where it was founded, with operating
subsidiaries and branches wherever it trades or employs people, each of which files where it is registered. Five
sources between them cover most of what is public.
\begin{itemize}
\item \textbf{Company registries.} The UK register is the most useful to this book's readers because it is free and
complete: every UK company and limited liability partnership files accounts, and anyone can download them. Registries
elsewhere differ in what they publish and whether they charge.
\item \textbf{SEC EDGAR.} US issuers file annual reports (Form 10-K) with audited statements; US broker-dealers file an
annual audited statement of financial condition, which is public, and periodic reports, which are not (below).
Chapter~1 used EDGAR's registration data to classify employers; this chapter reads the numbers.
\item \textbf{Annual reports of listed firms} elsewhere, which in the European Union are filed in a machine-readable
format (section~3).
\item \textbf{Prospectuses and offering documents}, when a firm lists or raises public debt. Private placements under
the US Rule~144A are offered to institutions only; what the press reports of their offering memoranda is context, not a
source.
\item \textbf{Divisional and Pillar~3 reports} of banks, which disclose a markets division's revenue (chapter~7) and,
uniquely in the industry, how many staff were paid more than a million euros.
\end{itemize}

\begin{dated}{2026-09}{The UK register's bulk accounts data}\label{dat:in:reading-a-trading-firms-accounts:registry}
Companies House publishes a free Accounts Data Product: daily and monthly ZIP files of the accounts filed electronically,
in \omterm{def:in:reading-a-trading-firms-accounts:ixbrl}{inline XBRL}, XBRL or zipped \omterm{def:in:reading-a-trading-firms-accounts:ixbrl}{inline XBRL}. It states that electronically filed accounts are ``about 75\% of the 2.2
million accounts we expect to be filed each year''; the rest, including many large companies' accounts, are filed as
paper or image documents, which the register also serves free, one filing at a time, but which carry no tags.
\end{dated}

\begin{definition}[FOCUS report]\label{def:in:reading-a-trading-firms-accounts:focus}
A \emph{FOCUS report}\index{FOCUS report} (Financial and Operational Combined Uniform Single report) is the periodic
report of a US broker-dealer's financial condition, income and net capital that it files with its regulators on Form
X-17A-5, Part~II or Part~IIA, within 17 business days of each quarter end (monthly for some firms). The reports are
confidential; the annual audited statement of financial condition is filed as a separate, public document.
\end{definition}

The split matters for a reader. A US market maker's broker-dealer publishes a balance sheet once a year and nothing
about its revenue or its staff: chapter~2 read one, Citadel Securities LLC's, for its assets and member's capital, and
that is as far as the public document goes. Its income, and the net-capital computation behind the rule Book~16,
chapter~2, describes, stay in the confidential FOCUS filings. For a private US trading firm, then, the UK accounts of
its London entity are often the only public income statement there is.

\begin{definition}[Pillar~3 disclosure]\label{def:in:reading-a-trading-firms-accounts:pillar3}
A \emph{Pillar~3 disclosure}\index{Pillar 3 disclosure} is the public report that prudential rules require of banks and
large investment firms on their capital, risks and remuneration. Its remuneration tables give, for the staff whose work
has a material impact on the firm's risk, the fixed and variable pay awarded, split by business area, and the number of
individuals paid one million euros or more, in bands of half a million euros up to five million and of one million
above.
\end{definition}

\begin{dated}{2025-12}{One bank's high earners}\label{dat:in:reading-a-trading-firms-accounts:pillar3}
Barclays PLC, Pillar~3 Report 2025, Appendix~C: 717 identified staff were paid \euro1 million or more in 2025 (UK REM4),
312 of them between \euro1 million and \euro1.5 million and 150 between \euro1.5 million and \euro2 million; 39 were paid
\euro5 million or more. Its investment bank's identified staff, 812.3 on the report's count, were awarded
\pounds1\,263.3 million, of which \pounds732.8 million variable (UK REM5).
\end{dated}

These are the only public pay distributions in the industry that name an employer. Of the bank's 717 high earners,
64.4\% were paid less than \euro2 million; the investment bank's identified staff averaged \pounds1.56 million, 58.0\%
of it variable. They describe the best-paid few hundred people of one bank, selected by the rules for their risk, not
its markets division as a whole; chapter~14 uses such tables with the care that selection needs.

\section{What statutory accounts contain}

A set of full UK accounts has the same parts whatever the business: a strategic report, a directors' report (for a partnership, a members' report), the auditor's report, the primary statements and the notes.
For a reader interested in a firm as an employer, four things are in the primary statements and the notes by law.

\begin{definition}[Staff costs, average headcount]\label{def:in:reading-a-trading-firms-accounts:staff}
\emph{Staff costs}\index{staff costs} are the costs of all persons employed by the entity during the year, broken
down into wages and salaries, social security costs, and other pension costs. \emph{Average
headcount}\index{average headcount} is the average number of persons employed under contracts of service during the
year: the sum over the year's months of the number employed in each month, divided by the number of months.
\end{definition}

Both come from the Companies Act 2006, section~411, which requires a company that is not small to give the average
number by category (the directors choose the categories) and the total \omterm{def:in:reading-a-trading-firms-accounts:staff}{staff costs} in the three parts. Each word does
work. ``Employed under contracts of service'' excludes contractors, secondees from another group company and
partners. ``Average'' means that a firm that doubled its staff during the year reports a figure between the two ends.
\omterm{def:in:reading-a-trading-firms-accounts:staff}{Staff costs} are what the year's profit and loss account is charged, which for deferred pay is the part expensed this
year, not the awards made this year (chapter~13). And a group's accounts count the group as one company.

\begin{definition}[Consolidated accounts]\label{def:in:reading-a-trading-firms-accounts:consolidated}
\emph{Consolidated accounts}\index{consolidated accounts} (group accounts) present a parent and its subsidiaries as a
single economic entity: their revenues, costs, assets and liabilities are added line by line, and transactions and
balances between members of the group are removed. The parent's own accounts, often printed beside them, show the
parent alone.
\end{definition}

A partnership's accounts add a fourth item. A limited liability partnership (Book~16, chapter~2) has members, not
shareholders; its profit is theirs.

\begin{definition}[Members' remuneration]\label{def:in:reading-a-trading-firms-accounts:members}
\emph{Members' remuneration}\index{members' remuneration} is what the members of a limited liability partnership
receive for the year: salaried amounts and profit shares that the partnership agreement allocates automatically, which
the accounts charge as an expense, and any profit left for the members to divide at their discretion. The accounts show
the profit before members' remuneration and profit shares, the average number of members, and, when that profit exceeds
\pounds200\,000, the profit attributable to the member with the largest entitlement.
\end{definition}

Two consequences follow for per-head arithmetic. First, the people who are members are not in the \omterm{def:in:reading-a-trading-firms-accounts:staff}{staff costs} or the
\omterm{def:in:reading-a-trading-firms-accounts:staff}{average headcount}: an LLP whose senior people are members understates the pay of its best-paid people by construction.
Second, members may be companies. A member that is itself a group company receives profit on behalf of the group, and
dividing the \omterm{def:in:reading-a-trading-firms-accounts:members}{members' remuneration} by the number of members then describes nothing (section~4).

\section{Machine-readable filings: inline XBRL and XBRL}

\begin{definition}[Inline XBRL]\label{def:in:reading-a-trading-firms-accounts:ixbrl}
\emph{Inline XBRL}\index{inline XBRL} is a format for filed accounts in which one XHTML document is both the page that a
person reads and the data that a program reads: each reported number is wrapped in a tag that names a concept from a
published taxonomy, a context (the entity, the period, and any dimension such as a segment or category), a unit, a
scale (a power of ten), a sign and the display format of the number.
\end{definition}

XBRL itself, the older form, is a separate machine-readable file beside the human-readable document. The US has
required it of issuers since 2009 and now receives it inline in 10-K and 10-Q filings; the European Union's single
electronic format requires issuers on its regulated markets to tag their consolidated IFRS statements in \omterm{def:in:reading-a-trading-firms-accounts:ixbrl}{inline XBRL},
and the accounts filed electronically at the UK register arrive as \omterm{def:in:reading-a-trading-firms-accounts:ixbrl}{inline XBRL} or XBRL. \Cref{lst:in:reading-a-trading-firms-accounts:tag}
shows a tagged fact from Quadrature Capital's accounts, stripped of its styling.
\omcode{code/firm/filings/tests/excerpt_quadrature_turnover.xhtml}{2}{3}{One tagged number from filed accounts: the
concept, context, unit, scale, decimals and format are attributes of the tag.}\label{lst:in:reading-a-trading-firms-accounts:tag}
It says: the concept turnover in the UK taxonomy, for the period that context \texttt{c2} defines (1 February 2024 to
31 January 2025), in the unit that \texttt{u1} defines (pounds sterling), displayed with commas as thousands
separators, to be multiplied by $10^3$ (\cref{fig:in:reading-a-trading-firms-accounts:ixbrl}).

\begin{omfigure}
\begin{tikzpicture}[font=\footnotesize,box/.style={draw=omDef,rounded corners=2pt,align=left,inner sep=4pt,fill=omDef!6}]
\node[box,text width=5.6cm] (tag) at (0,0) {\textbf{tagged number on the page}\\ name: concept (turnover)\\
contextRef $\to$ context\\ unitRef $\to$ unit\\ scale, sign, format\\ shown text: the printed digits};
\node[box,text width=5.0cm] (ctx) at (6.6,1.3) {\textbf{context c2}\\ entity: company number\\
period: \mbox{2024-02-01} to \mbox{2025-01-31}\\ dimensions: none};
\node[box,text width=5.0cm] (unit) at (6.6,-0.45) {\textbf{unit u1}\\ measure: iso4217:GBP};
\node[box,text width=5.0cm,draw=omThm,fill=omThm!6] (fact) at (6.6,-2.0) {\textbf{fact}\\ turnover, GBP, year to
\mbox{31 January} 2025\\ value $1\,224\,817\times10^3$};
\draw[-{Stealth},omDef] (tag.east) -- ++(0.35,0) |- (ctx.west);
\draw[-{Stealth},omDef] (tag.east) -- ++(0.35,0) |- (unit.west);
\draw[-{Stealth},omThm] (tag.south) |- (fact.west);
\end{tikzpicture}

\omcaption{Anatomy of one inline-XBRL fact. The tag on the page points to a context (entity, period, dimensions) and a
unit defined once in the document's hidden header; scale, sign and format turn the shown text into the value. The
brackets a statement prints around an expense are presentation, not sign. Example: Quadrature Capital Limited's
turnover for the year to 31 January 2025, read with \texttt{firm.filings.parse\_ixbrl}.}\label{fig:in:reading-a-trading-firms-accounts:ixbrl}
\end{omfigure}

The parser this chapter builds collects the contexts and units first, then turns every tagged number into a fact
(\cref{lst:in:reading-a-trading-firms-accounts:number}). Four details catch a first attempt. The format attribute
decides which character is the decimal mark: a continental filing writes 1.234,5. A dash tagged with the zero-dash
format is zero, not missing. The sign attribute, not the brackets around the displayed number, gives the fact's sign;
expenses are usually tagged positive although they print in brackets. And a number may be tagged twice, once on the face
of the statement and once in a note, which the parser keeps once.

\omcode{code/firm/filings/firm_filings.py}{72}{86}{One displayed number to its value: format, scale and
sign.}\label{lst:in:reading-a-trading-firms-accounts:number}

Tags are only as good as the filer's tagging. Quadrature's 2025 accounts tag other operating income of \pounds17.751
million with a minus sign, although the statement shows it as income that raises operating profit. A parser reports
what the tag says; a reconciliation catches what the tag gets wrong. Turnover less administrative expenses plus other
income, with the tagged sign, gives \pounds521.1 million against the stated operating profit of \pounds556.6 million, a
gap of \pounds35.5 million, twice the mistagged amount. With the sign reversed the three lines add up. Every snapshot in
this chapter is checked this way before it is used.

US filings reach a reader pre-parsed. The SEC's company-facts interface returns, for one issuer, every XBRL fact it has
ever filed, as JSON, with no key needed; it asks only that a script declare who it is in its user agent and stay under
ten requests a second. Each period recurs in every later filing that reports it as a comparative: Virtu Financial's
revenue concept holds 147 facts for its handful of annual periods. The reader keeps annual 10-K periods only and, for
each, the latest filing's value, because a later filing may restate
(\cref{lst:in:reading-a-trading-firms-accounts:companyfacts}).

\omcode{code/firm/filings/firm_filings.py}{169}{187}{Annual values from SEC company facts: the latest filing wins for
each period.}\label{lst:in:reading-a-trading-firms-accounts:companyfacts}

\section{Pitfalls: entity against group, members against employees, currencies and one-offs}

A correct fact read from the wrong entity, or divided by the wrong people, is a wrong answer. Jane Street's UK entities
show most of the pitfalls in six years.

\begin{dated}{2025-12}{One firm's two UK entities}\label{dat:in:reading-a-trading-firms-accounts:janestreet}
\footnotesize
\begin{tabular}{@{}lrrrrrr@{}}
\toprule
& 2020 & 2021 & 2022 & 2023 & 2024 & 2025\\
\midrule
\multicolumn{7}{@{}l}{\emph{Jane Street Europe Limited, group (\$m; \omterm{def:in:reading-a-trading-firms-accounts:staff}{average headcount})}}\\
revenues & 1\,567.0 & 222.5 & 435.0 & 383.4 & 995.8 & 643.0\\
profit for the year & 979.1 & $-69.1$ & 324.2 & 286.5 & 654.9 & 437.4\\
\omterm{def:in:reading-a-trading-firms-accounts:staff}{average headcount} & 274 & 335 & 2 & 3 & 4 & 5\\
\midrule
\multicolumn{7}{@{}l}{\emph{Jane Street UK Partnership LLP (\$m; 2022 is the 14 months from 25 October 2021)}}\\
revenues & & & 1\,230.8 & 1\,286.1 & 2\,271.1 & 3\,875.9\\
\omterm{def:in:reading-a-trading-firms-accounts:staff}{staff costs} & & & 434.9 & 494.2 & 860.1 & 1\,664.5\\
profit before members' shares & & & 667.5 & 583.5 & 1\,171.0 & 1\,855.4\\
\omterm{def:in:reading-a-trading-firms-accounts:staff}{average headcount} & & & 483 & 636 & 688 & 790\\
average members & & & 4 & 7 & 6 & 8\\
\bottomrule
\end{tabular}

\smallskip\noindent From the two entities' accounts at Companies House, 2021 to 2025, each figure with its page in
\texttt{data/industry/filings\_snapshot/jane\_street\_uk.csv}. On 1 August 2022 the company's employee contracts were
transferred to the partnership; from 1 January 2022 the company no longer received an allocation of employee costs.
\end{dated}

\textbf{Entity against group, and transfer pricing.} The company's revenue is what the group's transfer-pricing
arrangements (Book~16, chapter~15) assign to it for the trading it does, not what its traders earn in markets. In 2021
its \omterm{def:in:reading-a-trading-firms-accounts:staff}{staff costs}, \$302.4 million, exceeded its revenue of \$222.5 million, and it made a loss; in 2020 the same entity
had earned \$1\,567.0 million. The partnership's revenue is a charge to the group: in 2025 it billed \$2\,206.6 million to
its parent and \$1\,669.3 million to other affiliates, for activities performed on behalf of the group. Neither entity's
revenue is the London business's trading revenue, and the two cannot be added, because what one charges another group
company is, in the group's \omterm{def:in:reading-a-trading-firms-accounts:consolidated}{consolidated accounts}, removed. The \omterm{def:in:reading-a-trading-firms-accounts:consolidated}{consolidated accounts} that would settle the question
are those of the ultimate parent, Jane Street Group, LLC, in the United States, and are not among the UK filings.

\textbf{Where the people went.} \omterm{def:in:reading-a-trading-firms-accounts:staff}{Average headcount} fell from 335 to 2 in the company and appeared as 483 in the
partnership (\cref{fig:in:reading-a-trading-firms-accounts:jsheadcount}). The company's note counts the
individuals whose cost was allocated to it, and from 1 January 2022 almost none was. A reader with only the company's
2022 accounts would conclude that the London business had closed. With both, the staff cost per head is \$1.29 million in
2020 and \$0.90 million in 2021 (the company), then \$0.78 million, \$1.25 million and \$2.11 million in 2023--2025 (the
partnership).

\begin{omfigure}
\begin{tikzpicture}
\begin{axis}[width=10.5cm,height=5.4cm,ybar stacked,bar width=14pt,xmin=2019.4,xmax=2025.6,ymin=0,ymax=900,
  xtick={2020,2021,2022,2023,2024,2025},xticklabel style={/pgf/number format/1000 sep={}},
  ylabel={\small average headcount},tick label style={font=\footnotesize},grid=major,grid style={gray!20},
  table/col sep=comma,legend style={legend cell align=left,font=\scriptsize,at={(0.5,-0.2)},anchor=north,
  /tikz/every even column/.append style={column sep=8pt}},legend columns=2]
\addplot[fill=omDef!55,draw=omDef] table[x=year,y=group]{figdata/industry/11-reading-a-trading-firms-accounts/jane_street_headcount.csv};
\addplot[fill=omThm!55,draw=omThm] table[x=year,y=llp]{figdata/industry/11-reading-a-trading-firms-accounts/jane_street_headcount.csv};
\legend{company (group accounts),partnership (LLP)}
\end{axis}
\end{tikzpicture}

\omcaption{One firm's London headcount moves between its two UK entities. The partnership's first accounts cover the 14
months from 25 October 2021 to 31 December 2022 and are shown at 2022. Data: the entities' accounts at Companies House,
through \texttt{in\_accounts.jane\_street}.}\label{fig:in:reading-a-trading-firms-accounts:jsheadcount}
\end{omfigure}

\textbf{Periods.} The partnership's first accounts run 14 months. Its 2022 staff cost per head, \$0.90 million, is over
14 months; over twelve it is \$0.77 million. Annualise before comparing, and note that the first period of a new entity
also has a partial workforce.

\textbf{Members against employees.} In 2025 the partnership's profit before members' shares, \$1\,855.4 million, was
1.11 times its \omterm{def:in:reading-a-trading-firms-accounts:staff}{staff costs} and \$2.35 million per employee; it had eight members on average, ``a mix of corporate entities
and natural persons''. Dividing by eight gives \$232 million a member, a meaningless number: the member with the largest
entitlement, a corporate member, received 84.3\% of the profit (90.8\% in 2024). The profit leaves the partnership for the
group, not for eight people.

\textbf{An entity with no employees.} Squarepoint Capital LLP, a London sub-investment manager for its group's funds,
has corporate members only and no employees: the people who perform its regulated functions, 55 on average in 2025, are
seconded from an affiliate, and it pays a service charge for them, \pounds36.2 million in 2025. That is
\pounds659\,000 per secondee (\pounds738\,000 in 2024), but it is a charge between group companies, not pay, and
nothing in the accounts says what it contains. Its revenue halved from \pounds106.8 million to \pounds50.6 million in
2025 and its profit before \omterm{def:in:reading-a-trading-firms-accounts:members}{members' remuneration} fell from \pounds33.6 million to \pounds1.2 million; the members'
report attributes the fall in revenue to lower operating expenses and lower group income, which is what a service
entity's revenue follows.

\textbf{Currencies.} The two Jane Street entities report in dollars, Quadrature and Squarepoint in pounds. A comparison
converts at the average rate of the period (chapter~7's ECB rates), and for a year that ends in January, the calendar
year that holds eleven of its months.

\textbf{One-offs and what sits below operating profit.} For its fiscal years 2020 to 2023 Quadrature's profit before
tax exceeded its operating profit by \pounds234 million, \pounds250 million, \pounds351 million and \pounds179 million,
mostly fair-value gains on investments. Its revenue comes from managing funds: ``A portion of revenue is based on a fixed
fee structure, and the rest is calculated by reference to the performance of the funds under management.'' A margin
computed on operating profit and one computed on profit before tax describe different businesses. Distributions have
the same trap: in fiscal 2024 the statement of changes in equity records dividends of \pounds1\,329.4 million, while the
cash-flow statement records \pounds82.9 million of dividends paid in cash.

\section{Per-head measures and their caveats}

With the pitfalls in mind, six years of one entity whose employees, costs and revenue sit in the same accounts read
cleanly. Quadrature Capital Limited is such an entity: a company, not a partnership, filing full accounts with the
average number of employees and \omterm{def:in:reading-a-trading-firms-accounts:staff}{staff costs} in its notes.

\begin{dated}{2025-01}{Six years of one London investment manager}\label{dat:in:reading-a-trading-firms-accounts:quadrature}
\footnotesize
\begin{tabular}{@{}lrrrrrrrr@{}}
\toprule
fiscal year & turnover & staff & employ- & staff cost & turnover & staff/ & profit & dividends\\
to January & \pounds m & costs \pounds m & ees & per head & per head & turnover & after tax & \pounds m\\
\midrule
2020 & 140.6 & 116.7 & 68 & 1.72 & 2.07 & 83.0\% & 244.2 & 4.0\\
2021 & 269.8 & 182.9 & 83 & 2.20 & 3.25 & 67.8\% & 262.7 & 1.0\\
2022 & 778.6 & 409.0 & 98 & 4.17 & 7.94 & 52.5\% & 525.6 & 0.0\\
2023 & 664.2 & 394.7 & 113 & 3.49 & 5.88 & 59.4\% & 225.6 & 0.0\\
2024 & 588.2 & 408.0 & 143 & 2.85 & 4.11 & 69.4\% & 49.6 & 1\,329.4\\
2025 & 1\,224.8 & 556.1 & 173 & 3.21 & 7.08 & 45.4\% & 410.7 & 360.0\\
\bottomrule
\end{tabular}

\smallskip\noindent Per-head figures in \pounds m. Fiscal 2024 and 2025 from the inline-XBRL accounts for the year to
January 2025; 2020--2023 transcribed from the image accounts for the years to January 2021 and 2023. Every figure with
its filing and page in \texttt{data/industry/filings\_snapshot/}.
\end{dated}

Staff cost per head ranged from \pounds1.72 million to \pounds4.17 million over the six years, and turnover per head from
\pounds2.07 million to \pounds7.94 million (\cref{fig:in:reading-a-trading-firms-accounts:quadrature,fig:in:reading-a-trading-firms-accounts:perhead}).
\omterm{def:in:reading-a-trading-firms-accounts:staff}{Staff costs} took between 45.4\% and 83.0\% of turnover: in the lean years they took most of it, and in the best years
less than half. The pattern is Book~16, chapter~1's compensation ratio at work, with variable pay absorbing part of the
swing in revenue; that it does so here is visible, how it does so is not.

\begin{omfigure}
\begin{tikzpicture}
\begin{axis}[width=10.5cm,height=5.4cm,ybar,area legend,bar width=6pt,xmin=2019.4,xmax=2025.6,ymin=0,ymax=1400,
  xtick={2020,2021,2022,2023,2024,2025},xticklabel style={/pgf/number format/1000 sep={}},
  xlabel={\small fiscal year to January},ylabel={\small \pounds{} million},tick label style={font=\footnotesize},
  grid=major,grid style={gray!20},table/col sep=comma,legend style={legend cell align=left,font=\scriptsize,at={(0.5,-0.3)},anchor=north,
  /tikz/every even column/.append style={column sep=8pt}},legend columns=3]
\addplot[fill=omDef!55,draw=omDef] table[x=fy,y=revenue]{figdata/industry/11-reading-a-trading-firms-accounts/quadrature_years.csv};
\addplot[fill=omThm!55,draw=omThm] table[x=fy,y=staff]{figdata/industry/11-reading-a-trading-firms-accounts/quadrature_years.csv};
\addplot[fill=gray!40,draw=gray] table[x=fy,y=profit]{figdata/industry/11-reading-a-trading-firms-accounts/quadrature_years.csv};
\legend{turnover,staff costs,profit after tax}
\end{axis}
\end{tikzpicture}

\omcaption{Six fiscal years of one London investment manager: turnover, \omterm{def:in:reading-a-trading-firms-accounts:staff}{staff costs} and profit after tax. Profit after
tax exceeds turnover less costs in 2020--2023 because of fair-value gains on investments below operating profit. Data:
the company's accounts at Companies House, through \texttt{in\_accounts.quadrature}.}\label{fig:in:reading-a-trading-firms-accounts:quadrature}
\end{omfigure}

\begin{omfigure}
\begin{tikzpicture}
\begin{axis}[width=10.5cm,height=5.2cm,xmin=2019.6,xmax=2025.4,ymin=0,ymax=9,
  xtick={2020,2021,2022,2023,2024,2025},xticklabel style={/pgf/number format/1000 sep={}},
  xlabel={\small fiscal year to January},ylabel={\small \pounds{} million per employee},tick label style={font=\footnotesize},
  grid=major,grid style={gray!20},table/col sep=comma,legend style={legend cell align=left,font=\scriptsize,at={(0.5,-0.3)},anchor=north,
  /tikz/every even column/.append style={column sep=8pt}},legend columns=2]
\addplot[omDef,thick,mark=*] table[x=fy,y=revenue_per_head]{figdata/industry/11-reading-a-trading-firms-accounts/quadrature_perhead.csv};
\addplot[omThm,thick,mark=square*] table[x=fy,y=staff_per_head]{figdata/industry/11-reading-a-trading-firms-accounts/quadrature_perhead.csv};
\legend{turnover per employee,staff costs per employee}
\end{axis}
\end{tikzpicture}

\omcaption{The same entity per average employee. Staff cost per head moves with turnover per head but by less: from
\pounds1.72 million to \pounds4.17 million while turnover per head moved from \pounds2.07 million to \pounds7.94
million. Data: as \cref{fig:in:reading-a-trading-firms-accounts:quadrature}.}\label{fig:in:reading-a-trading-firms-accounts:perhead}
\end{omfigure}

Six caveats go with every per-head figure read from accounts.
\begin{enumerate}
\item \textbf{A mean over a skewed distribution.} In a firm whose pay is mostly variable, a few people take a large
share; the typical employee earns less than the mean, and chapter~14 gives the medians that other sources allow.
\item \textbf{The numerator includes the employer's costs.} \omterm{def:in:reading-a-trading-firms-accounts:staff}{Staff costs} include the employer's social security costs
(in Quadrature's 2025 accounts, \pounds74.4 million of \pounds556.1 million) and pension contributions, which the
employee never sees as pay.
\item \textbf{The denominator is an average.} A firm growing from 143 to 173 employees in a year spreads the year's
costs over people who were there for part of it.
\item \textbf{Timing.} Deferred awards reach the profit and loss account over their vesting period, so a year's \omterm{def:in:reading-a-trading-firms-accounts:staff}{staff
costs} mix this year's pay with earlier years' awards (chapter~13).
\item \textbf{Who is missing.} Partners and members, contractors and secondees are outside both halves of the fraction.
\item \textbf{Denominators of ratios.} Virtu Financial's employee compensation took 12.1\% to 17.2\% of its reported
revenues in 2020--2025, and 15.8\% to 30.3\% of its trading income alone. A market maker's reported revenues are
before its volume-driven costs (Book~16, chapter~1), a fund manager's turnover is fees; a ratio built on one line is not
comparable with a ratio built on another. Divided by the roughly 1\,027 employees its 10-K reports at the start of 2026,
its 2025 compensation was about \$0.51 million a head, from a count at a date rather than an average.
\end{enumerate}

Converted at 2024's average rate of \$1.2785 to the pound, Quadrature's staff cost per head for the year to January
2025 is \$4.11 million, against \$1.25 million and \$2.11 million in Jane Street's partnership in 2024 and 2025. The
comparison is legitimate as arithmetic and weak as evidence: a small investment manager's employees against a large
market maker's London operation, whose members' profit shares, larger than its \omterm{def:in:reading-a-trading-firms-accounts:staff}{staff costs}, are outside the fraction.

\section{Tutorial: five years at Companies House}\label{tut:in:reading-a-trading-firms-accounts:tut}

\begin{tutorial}
\textbf{Goal.} Build a six-year table of one firm's accounts from its filings, check it, and compute its per-head
measures. \textbf{End state:} the table in the dated box above and
\cref{fig:in:reading-a-trading-firms-accounts:quadrature,fig:in:reading-a-trading-firms-accounts:perhead}.
\begin{tutsteps}
\item \textbf{Fetch.} The chapter's script \texttt{in\_fetch.py}, run with a scratch directory as its argument, downloads the
inline-XBRL accounts and SEC company facts into a scratch directory, which is never committed, and writes
\texttt{quadrature\_ixbrl.csv} and \texttt{virtu.csv} in \texttt{data/industry/filings\_snapshot/}. The tests never touch
the network: they read the snapshot, and a synthetic inline-XBRL document in \texttt{code/firm/filings/tests/}.
\item \textbf{Parse.} \texttt{firm.filings.parse\_ixbrl(text, entity, source)} returns the facts, with the taxonomy
concept kept in each fact's note and the book's short name (\texttt{Revenue}, \texttt{StaffCosts},
\texttt{AverageEmployees}) as its concept; concepts about a named director are dropped.
\item \textbf{Transcribe.} The older accounts are images. Their lines go by hand into
\texttt{quadrature\_pdf.csv}, one row per fact with its filing's URL and page; a second reading checks every row.
\item \textbf{Check.} \texttt{Snapshot.load(\dots).problems()} lists facts without provenance and periods stated twice
with different values; \texttt{reconcile} tests that operating profit equals its parts
(\cref{lst:in:reading-a-trading-firms-accounts:reconcile}) and finds the mistagged sign of fiscal 2025.
\omcode{code/firm/filings/firm_filings.py}{248}{258}{The snapshot checks its own arithmetic: a stated total against
the sum of its parts.}\label{lst:in:reading-a-trading-firms-accounts:reconcile}
\item \textbf{Measure.} \texttt{in\_accounts.quadrature()} computes the table; \texttt{firm.filings.to\_statement} maps
each year onto Book~16's \texttt{Statement} for the ratios and the cost structure of \texttt{firm.firmecon}.
\end{tutsteps}
\end{tutorial}

\textbf{What to change next.} Add the two missing Quadrature years' own filings (fiscal 2022 and 2024) and check that
their figures equal the comparatives used here; add a firm whose accounts are filed as \omterm{def:in:reading-a-trading-firms-accounts:ixbrl}{inline XBRL} for all five years
and write a snapshot without transcription.

\section{Build: the filings reader}\label{bld:in:reading-a-trading-firms-accounts:filings}

\begin{build}
\bfield{Purpose} Read filed accounts into facts with provenance, for this chapter's firm, chapter~12's industry table
and chapter~14's pay sources.
\bfield{Interface} \texttt{firm.filings}: \texttt{Fact}; \texttt{number(text, fmt, scale, sign)};
\texttt{parse\_ixbrl(text, entity, source, keep\_private)}; \texttt{companyfacts(doc, concepts, entity, source)};
\texttt{Snapshot} with \texttt{load}, \texttt{save}, \texttt{get}, \texttt{series}, \texttt{ends}, \texttt{problems},
\texttt{reconcile}; \texttt{months}; \texttt{per\_head}; \texttt{to\_statement(snap, entity, end, mapping, headcount)}.
Python standard library only (the mapping reuses \texttt{firm.firmecon}).
\bfield{Rules} Every fact carries a source and a locator (page, context or accession); the tag's sign is the value's
sign; identical repeats are kept once; dimensional facts are kept but ignored by lookups; for company facts, annual
10-K periods only and the latest filing per period; concepts about named individuals are dropped unless explicitly
requested.
\bfield{Acceptance tests} \texttt{code/firm/filings/tests/}: on a synthetic document, formats, scales and signs give
the right values, the duplicate and the nil fact are dropped, the dimensional fact keeps its member, the private
concept is dropped, and the reconciliation flags the mistagged sign; on synthetic company facts, the restated value
and not the first one is kept and quarterly periods are ignored.
\bfield{Stretch} Read ESEF packages (a ZIP with the report and its taxonomy extension); resolve concepts through the
taxonomy's labels instead of a fixed alias table; parse the Companies House bulk files in a stream.
\end{build}

\begin{omsources}
\item Companies Act 2006, section~411; The Large and Medium-sized Limited Liability Partnerships (Accounts)
Regulations 2008, Schedule~1.
\item Companies House, Free Accounts Data Product; the accounts of Quadrature Capital Limited (years to January 2021,
2023 and 2025), Jane Street Europe Limited (2021, 2022, 2023 and 2025), Jane Street UK Partnership LLP (2023 and 2025)
and Squarepoint Capital LLP (2025).
\item SEC, Accessing EDGAR Data and EDGAR Application Programming Interfaces; XBRL company facts of Virtu Financial,
Inc.; 17 CFR 240.17a-5; FINRA, eFOCUS.
\item ESMA, Electronic Reporting (ESEF).
\item Regulation (EU) No~575/2013, Article~450; Barclays PLC, Pillar~3 Report 2025.
\end{omsources}

\section{Exercises}

\begin{exercise}[$\star$]\label{exo:in:reading-a-trading-firms-accounts:1}
From the facts of Quadrature's accounts for the year to January 2025 (turnover \pounds1\,224.8 million, \omterm{def:in:reading-a-trading-firms-accounts:staff}{staff costs}
\pounds556.1 million, 173 average employees), compute staff cost per head, turnover per head and \omterm{def:in:reading-a-trading-firms-accounts:staff}{staff costs} as a share of
turnover.
\end{exercise}

\begin{exercise}[$\star$]\label{exo:in:reading-a-trading-firms-accounts:2}
Jane Street's partnership reported \omterm{def:in:reading-a-trading-firms-accounts:staff}{staff costs} of \$434.9 million and an average of 483 employees for its first
period, 25 October 2021 to 31 December 2022. Compute staff cost per head for the period and for a twelve-month year.
\end{exercise}

\begin{exercise}[$\star$]\label{exo:in:reading-a-trading-firms-accounts:3}
From the bank's Pillar~3 bands in the dated box, how many of its staff were paid \euro1 million or more, and what share of
them less than \euro2 million?
\end{exercise}

\begin{exercise}[$\star\star$]\label{exo:in:reading-a-trading-firms-accounts:4}
Quadrature's 2025 facts give turnover \pounds1\,224.817 million, administrative expenses \pounds686.012 million, other
operating income tagged at $-$\pounds17.751 million and operating profit \pounds556.556 million. Reconcile, find the
error, and say how a program should report it.
\end{exercise}

\begin{exercise}[$\star\star$]\label{exo:in:reading-a-trading-firms-accounts:5}
Convert Quadrature's staff cost per head for the year to January 2025 to dollars at 2024's average rate and compare it
with Jane Street's partnership in 2024 and 2025. Give two reasons the comparison is weak.
\end{exercise}

\begin{exercise}[$\star\star$]\label{exo:in:reading-a-trading-firms-accounts:6}
Why does dividing Jane Street's partnership's 2025 profit before members' shares by its average number of members
mislead? What share went to the largest member?
\end{exercise}

\begin{exercise}[$\star\star\star$]\label{exo:in:reading-a-trading-firms-accounts:7}
\emph{Coding.} Load \texttt{virtu.csv} with \texttt{firm.filings} and compute employee compensation as a share of
revenues and of trading income for 2020--2025. Which denominator moves more, and why?
\end{exercise}

\begin{exercise}[$\star\star\star$]\label{exo:in:reading-a-trading-firms-accounts:8}
\emph{Find the flaw.} ``Jane Street Europe Limited earned \$643 million of revenue in 2025 with five employees: \$129
million per employee, the most productive trading staff in London.''
\end{exercise}

\section{Problem: Five Years at Companies House}

\begin{problem}[{Weekend problem --- five years at Companies House}]\label{pb:in:reading-a-trading-firms-accounts:1}
A candidate with an offer from a London firm wants to know what its accounts say about pay, and how far to trust it.
She chooses a company that files full accounts and has six years on the register.

\medskip\noindent\textbf{Part I --- Finding the filings.}
\begin{enumerate}
\item Define a \omterm{def:in:reading-a-trading-firms-accounts:registry}{company registry} and \omterm{def:in:reading-a-trading-firms-accounts:registry}{statutory accounts}.
\item Name five kinds of public source for a trading firm's finances.
\item What share of UK accounts does the register receive in machine-readable form, and what are the others?
\item Define a \omterm{def:in:reading-a-trading-firms-accounts:focus}{FOCUS report}. Which parts of a US broker-dealer's filings are public?
\item Define a \omterm{def:in:reading-a-trading-firms-accounts:pillar3}{Pillar~3 disclosure}. What does it say about pay that no other filing does?
\end{enumerate}

\medskip\noindent\textbf{Part II --- Reading them.}
\begin{enumerate}[resume]
\item Define \omterm{def:in:reading-a-trading-firms-accounts:staff}{staff costs} and \omterm{def:in:reading-a-trading-firms-accounts:staff}{average headcount}, and give the law's three parts of \omterm{def:in:reading-a-trading-firms-accounts:staff}{staff costs}.
\item Who is not in the \omterm{def:in:reading-a-trading-firms-accounts:staff}{average headcount}?
\item Define \omterm{def:in:reading-a-trading-firms-accounts:consolidated}{consolidated accounts} and \omterm{def:in:reading-a-trading-firms-accounts:members}{members' remuneration}.
\item Define \omterm{def:in:reading-a-trading-firms-accounts:ixbrl}{inline XBRL}. What do the context, unit, scale, sign and format of a tag give?
\item Why does the company-facts reader keep only the latest filing for each period?
\end{enumerate}

\medskip\noindent\textbf{Part III --- The six years.}
\begin{enumerate}[resume]
\item Give the firm's turnover, \omterm{def:in:reading-a-trading-firms-accounts:staff}{staff costs} and average employees for fiscal 2020 to 2025.
\item Compute staff cost per head and turnover per head for each year.
\item Compute \omterm{def:in:reading-a-trading-firms-accounts:staff}{staff costs} as a share of turnover. When is it highest, and why?
\item Reconcile fiscal 2025's operating profit and explain the discrepancy.
\item Why did profit after tax exceed operating profit in fiscal 2020 to 2023?
\end{enumerate}

\medskip\noindent\textbf{Part IV --- The verdict.}
\begin{enumerate}[resume]
\item Over six years, what share of profit after tax was paid out as dividends?
\item State the \emph{named result}: the ranges of staff cost per head and turnover per head over the six years, and
the share of profit after tax distributed.
\item Give three reasons the mean staff cost per head overstates what she will be paid.
\item Name two pitfalls that would have misled her had she chosen a partnership or a group entity instead.
\item In two sentences, what do the accounts tell her, and what do they not?
\end{enumerate}
\end{problem}

\section{Interview questions}

\begin{interviewq}[$\star$ \iqroles{researcher, developer}]\label{iq:in:reading-a-trading-firms-accounts:1}
A filing shows an expense as (900). What value should a parser store, and where does the sign come from?
\end{interviewq}

\begin{interviewq}[$\star$ \iqroles{bank, risk}]\label{iq:in:reading-a-trading-firms-accounts:2}
What is the difference between a bank's \omterm{def:in:reading-a-trading-firms-accounts:consolidated}{consolidated accounts} and its parent company's accounts, and which would you
use to compare it with another bank?
\end{interviewq}

\begin{interviewq}[$\star\star$ \iqroles{developer, mle}]\label{iq:in:reading-a-trading-firms-accounts:3}
Design a pipeline that keeps a table of every UK trading firm's \omterm{def:in:reading-a-trading-firms-accounts:staff}{staff costs} and headcount up to date from the
register's daily files. What do you store, and how do you handle restatements?
\end{interviewq}

\begin{interviewq}[$\star\star$ \iqroles{researcher}]\label{iq:in:reading-a-trading-firms-accounts:4}
A firm's revenue per employee doubled in a year. List the explanations that are not ``the staff became more
productive''.
\end{interviewq}

\begin{interviewq}[$\star\star$ \iqroles{trader, bank}]\label{iq:in:reading-a-trading-firms-accounts:5}
Why might a trading firm book its London revenue in one UK entity and its London staff in another?
\end{interviewq}

\begin{interviewq}[$\star\star\star$ \iqroles{researcher, risk}]\label{iq:in:reading-a-trading-firms-accounts:6}
Given six years of a firm's revenue and \omterm{def:in:reading-a-trading-firms-accounts:staff}{staff costs}, how would you estimate how much of its pay is variable, and what
would make the estimate wrong?
\end{interviewq}
